\documentclass[11pt]{article}
\usepackage[T1]{fontenc}
\usepackage[margin=1in]{geometry}
\usepackage{amsmath,amssymb}
\usepackage[hidelinks]{hyperref}

\title{Literature Status: One-Factor Projective Tensor Octahedrality}
\author{Automatic Math Discovery Run \texttt{fa\_banach\_001}}
\date{June 28, 2026}

\begin{document}
\maketitle

\section*{Status}

This is a \textbf{literature\_already\_answered} packet.  It records that a
later paper explicitly answers one of the open questions extracted from
arXiv:1609.02062.  It is not a new proof packet.

\section*{Source Question}

The source paper is:

\begin{quote}
Juhan Langemets, Vegard Lima, and Abraham Rueda Zoca,
\emph{Octahedral norms in tensor products of Banach spaces},
arXiv:1609.02062.
\end{quote}

Near the end of the paper, after proving a partial positive result for
projective tensor products with a non-reflexive \(L\)-embedded factor, the
authors write that for general Banach spaces \(X\) and \(Y\) it remains open
whether
\[
  X \widehat\otimes_\pi Y
\]
has an octahedral norm whenever \(X\) and/or \(Y\) has an octahedral norm.
The source also describes the partial result as a partial positive answer to
\cite[Question~4.4]{LLR2017}, where the one-factor projective preservation
problem is asked.

\section*{Supporting Answer}

The supporting paper is:

\begin{quote}
Abraham Rueda Zoca,
\emph{Stability results of octahedrality in tensor product spaces},
arXiv:1708.07300.
\end{quote}

The abstract of the supporting paper states that there exists a finite
dimensional complex Banach space \(X\) such that
\[
  L_1^\mathbb C([0,1]) \widehat\otimes_\varepsilon X
\]
fails the strong diameter two property and
\[
  L_\infty^\mathbb C([0,1]) \widehat\otimes_\pi X^*
\]
fails to have an octahedral norm.  It further states that this proves that
octahedrality is not automatically inherited from one factor by taking
projective tensor product, answering \cite[Question~4.4]{LLR2017}.

The paper proves this through Theorem~2.1 and Theorem~2.2.  In particular,
Theorem~2.2 gives a two-dimensional complex Banach space \(Y\) such that
\[
  L_\infty^\mathbb C([0,1]) \widehat\otimes_\pi Y
\]
does not have an octahedral norm.  The following remark identifies why this is
a counterexample to the one-factor problem: \(L_\infty^\mathbb C([0,1])\) has
the Daugavet property, hence has an octahedral norm, but the projective tensor
product above is not octahedral.

\section*{Identification}

The final open question in arXiv:1609.02062 asks whether octahedrality of one
factor forces octahedrality of the projective tensor product in general.  The
supporting paper constructs a projective tensor product
\[
  L_\infty^\mathbb C([0,1]) \widehat\otimes_\pi Y
\]
where the first factor is octahedral but the tensor product is not.  This is a
negative answer to the same one-factor projective tensor product question.

The bibliographic identification is explicit: arXiv:1708.07300 cites
\cite{LLR2017} as Question 4.4 and cites the arXiv:1609.02062 source paper as
\cite{LLR2017QJM}.  Thus this packet is best classified as
\texttt{literature\_already\_answered}, not as a new counterexample.

\section*{Scope Limitations}

This packet clears only the projective tensor product octahedrality question.
It does not claim a new proof of the source paper's partial positive theorem,
and it does not address unrelated diameter-two or injective tensor product
stability questions except insofar as they are used in the supporting paper's
duality argument.

\section*{Search Evidence}

Local cheap indexes for arXiv:1609.02062, the title, octahedral norms,
projective tensor products, and Question 4.4 had no prior exact packet at
reservation time.  A targeted literature search found arXiv:1708.07300, whose
abstract and post-Theorem~2.2 remark explicitly state the negative answer to
Question 4.4.

\begin{thebibliography}{9}

\bibitem{LLR2017}
J. Langemets, V. Lima, and A. Rueda Zoca,
\emph{Almost square and octahedral norms in tensor products of Banach spaces},
Rev. R. Acad. Cienc. Exactas Fis. Nat. Ser. A Mat. \textbf{111} (2017),
841--853.

\bibitem{LLR2017QJM}
J. Langemets, V. Lima, and A. Rueda Zoca,
\emph{Octahedral norms in tensor products of Banach spaces},
Q. J. Math. \textbf{68} (2017), no. 4, 1247--1260; arXiv:1609.02062.

\bibitem{RZ1708}
A. Rueda Zoca,
\emph{Stability results of octahedrality in tensor product spaces},
arXiv:1708.07300.

\end{thebibliography}

\end{document}
